\section{How to Use This Expansion}

This expansion is designed for players who want to build a Runekeeper who fights, endures, and carries their covenant in flesh and steel. It is also for Game Masters who want to introduce Runic Warriors into their campaigns as allies, antagonists, or complications.

\subsection{For Players}

You are building a warrior who is also a Runekeeper. You are not a scholar. You are not a courtier. You are the one who stands at the front, who holds the line, who charges when others hesitate.

\subsubsection{What You Need}

Before using this expansion, you should be familiar with the core Runekeeper rules from the \textit{Fate's Edge Comprehensive Resource Guide}. This expansion does not replace those rules. It builds on them.

You will still:
\begin{itemize}
\item Choose a Patron and track Obligation.
\item Invoke Rites and mark costs.
\item Manage your Codex and Thiasos.
\end{itemize}

But your Codex will not be a book. Your Thiasos may be a warhorse, a spirit-hound, or a sentient blade. Your Rites will be carved, engraved, or tattooed. And your Obligation will be written on your body.

\subsubsection{How to Build a Runic Warrior}

\begin{enumerate}
\item \textbf{Choose a Patron.} Not every Patron is suited to a Runic Warrior. This expansion identifies twenty-eight Patrons that welcome physical covenants, organized into thematic families: The Vowed, The Threshold, The Wild, The Deep, The Shadow, The Web, The Forge, The Grave, The Fracture, and the Lethai Ghe'hai.

\item \textbf{Choose your Codex form.} Your Codex is not a book. It is an engraved gorget, a carved staff, a tattooed body, a sentient blade. Each form has advantages and risks. A gorget grants armor. A staff serves as a weapon. Tattoos cannot be stolen. A sentient blade offers counsel but may have opinions.

\item \textbf{Choose your Thiasos.} Your Thiasos is not a pet. It is a warhorse, a spirit-hound, a hollowed bird, an ancestral shade, or a clockwork construct. It fights beside you, witnesses your Rites, and shares your burdens. Optional rules allow for sentient Thiasos—weapons or objects that speak and remember.

\item \textbf{Select Runic Warrior talents.} This expansion introduces new talents designed for physical resilience, close combat, and the conversion of Fatigue into advantage. Foundation talents (2-4 XP) are available at Tier I. Advanced talents (6-8 XP) unlock at Tier II. Prestige talents (10-12 XP) are for veteran warriors.

\item \textbf{Learn new Rites.} Each Patron family includes new Rites tailored to Runic Warriors. Low Rites (4-5 XP) are quick and scene-limited. Standard Rites (7-9 XP) offer lasting effects. High Rites (12-14 XP) are powerful, costly, and often require a witness or threshold.
\end{enumerate}

\subsubsection{Build Archetypes}

This expansion includes five sample archetypes to guide your character creation:

\begin{itemize}
\item \textbf{Dawn Warrior (Oath of Flame \& Light)} — A crusader who smites foes with radiant strikes and swears vows before battle.
\item \textbf{Bloodsworn (High Mother)} — A Daughter of the Cord who stores Fatigue as visible knots in the flesh.
\item \textbf{Ghe'hai Operative (Lethai Syncretic)} — An elite warrior-assassin who wears a silk mask and binds foes in invisible threads.
\item \textbf{Beast Warrior (Savage Heart)} — A hunter who wears the skin of a killed predator and fights with claw and fang.
\item \textbf{Fractal Warrior (The Ninth)} — A warrior who fights with impossible angles and precognition, their Codex a broken mirror.
\end{itemize}

\subsection{For Game Masters}

You are introducing a new kind of Runekeeper into your world. These are not wizards in armor. They are warriors who have made covenants with forces that demand physical proof of loyalty.

\subsubsection{Integration into Existing Campaigns}

Runic Warriors can appear anywhere, but their form varies by region.

\begin{itemize}
\item In \textbf{Aeler}, they are Cairn-Knights or Artificers, their Rites carved into burial stones or toolbelts.
\item In \textbf{Sekogo}, they are Beast Warriors, wearing predator skins and hunting as sacrament.
\item In \textbf{Ashaan}, they are Shadow Warriors, moving unseen, their Thiasos a raven that carries whispers.
\item In \textbf{Lethai}, they are Ghe'hai operatives—silks, steel, and silence.
\item In \textbf{Viterra} and \textbf{Ecktoria}, they are Justicars, Templars, and Witch-Hunters, their Codex a shield or a censer.
\end{itemize}

Use the faction relationship matrix in the appendices to determine how local powers view Runic Warriors. The Chain-Lanterns may be suspicious of a Runekeeper whose Codex cannot be confiscated. The Daughters of the Cord may see a Bloodsworn as a sister—or a rival.

\subsubsection{Optional Rules}

This expansion includes optional rules for sentient Thiasos: weapons, shields, or objects that speak and remember. A sentient Thiasos cannot act independently, but it can serve as a witness, offer counsel, and carry Rites engraved on its surface.

Use sentient Thiasos as:
\begin{itemize}
\item Allies that provide information or warnings.
\item Complications that have their own agendas and opinions.
\item Narrative devices that carry secrets or memories from previous owners.
\end{itemize}

\subsubsection{Sample NPCs}

The GM guidance chapter includes sample Runic Warrior NPCs for each Patron family. Use them as antagonists, allies, or neutral parties.

\subsection{Cultural Integration}

Runic Warriors do not belong to a single culture. They emerge wherever warriors make covenants. This expansion provides tools to build a Runic Warrior from any people, under any Patron.

The thread held. The water carried. That is not a Rite. That is grace.
